\documentclass[aspectratio=169,11pt]{beamer}
\usepackage[utf8]{inputenc}
\usepackage[T1]{fontenc}
\usepackage[french]{babel}
\usepackage{helvet}
\renewcommand{\familydefault}{\sfdefault}
\usepackage{hyperref}
\definecolor{bleu}{HTML}{0000FF}
\definecolor{fond}{HTML}{F4F4EF}
\setbeamercolor{background canvas}{bg=fond}
\setbeamercolor{structure}{fg=bleu}
\setbeamercolor{title}{fg=bleu}
\setbeamercolor{frametitle}{fg=bleu,bg=fond}
\setbeamercolor{normal text}{fg=black}
\setbeamertemplate{navigation symbols}{}
\setbeamertemplate{itemize item}{\textcolor{bleu}{\tiny$\blacksquare$}}
\setbeamertemplate{footline}{\hfill\footnotesize\insertframenumber\hspace{1em}\vskip6pt}
\hypersetup{colorlinks=true,urlcolor=bleu,linkcolor=bleu}
\title{Prêts et programmes}
\subtitle{Fonds monétaire international}
\author{Veille institutionnelle}
\date{2026-11-03}

\begin{document}
\begin{frame}\titlepage\end{frame}

\begin{frame}{Prêts et programmes (1/2)}
\begin{itemize}
\item SBA et EFF pour les besoins de balance des paiements
\item RFI et RCF : financements d'urgence
\item FCL, PLL, SLL : lignes de précaution
\item PRGT : prêts concessionnels. RSF : résilience climatique
\end{itemize}
\end{frame}

\begin{frame}{Prêts et programmes (2/2)}
\begin{itemize}
\item Conditionnalité : actions préalables, critères de réalisation, repères structurels
\item Dette : cadre commun du G20, table ronde mondiale, analyse de viabilité
\end{itemize}
\end{frame}

\begin{frame}{Prêts et programmes : ressources à suivre}
\small
Veille hebdomadaire.
\begin{itemize}
\item \href{https://www.imf.org/en/About/Factsheets/Sheets/2023/IMF-Lending}{Prêts du FMI}\\{\footnotesize Fiche institutionnelle, publication : annuelle}
\item \href{https://www.imf.org/en/Countries/ResRep}{Programmes en cours}\\{\footnotesize Page pays, publication : mensuelle}
\item \href{https://www.imf.org/en/Topics/imf-and-debt}{Dette et FMI}\\{\footnotesize Page thématique, publication : trimestrielle}
\end{itemize}
\end{frame}

\end{document}
